\begin{figure}[tp]
\addtocounter{figure}{-1}
\centering
\begin{subfigure}{1.0\linewidth}
\centering
\videoframe{figures/video_frames/mona_yawn_lores_00001.jpg}
\videoframe{figures/video_frames/mona_yawn_lores_00006.jpg}
\videoframe{figures/video_frames/mona_yawn_lores_00011.jpg}
\videoframe{figures/video_frames/mona_yawn_lores_00017.jpg}
{\textbf{Animated from still image}
\ \\
\ \\}
\end{subfigure}
% \vspace{-4mm}


\begin{subfigure}{1.0\linewidth}
\centering
\videoframe{figures/video_frames/mona_yawn_snow_00001.jpg}
\videoframe{figures/video_frames/mona_yawn_snow_00006.jpg}
\videoframe{figures/video_frames/mona_yawn_snow_00011.jpg}
\videoframe{figures/video_frames/mona_yawn_snow_00017.jpg}
{\textbf{Stylized video} \\
\textbf{Prompt:} An oil painting of a snowman with a red hat opening their mouth to yawn
\ \\}
\end{subfigure}
\captionof{figure}{\textbf{Example of zero-shot video editing via task chaining} (text conditioned image-to-video and stylization) 
%-- the original painting is first animated via a text prompt and then stylized via another text prompt.
}
\label{fig:animate_stylize}
\vspace{-2mm}
\end{figure}


% %\begin{figure*}[h]
% \begin{figure}[t]
% % \addtocounter{figure}{-1} %Prevent subfigures increasing main figure count  % for some reason, this needed to be commented out to get this figure to have a different number than the preceeding figure.
% \centering
% % Arc shot
% \begin{subfigure}[b]{1.0\linewidth}
% %{0.5\textwidth}
% \centering
% \videoframe{figures/video_frames/train_original_00001.jpg}
% \videoframe{figures/video_frames/train_original_00006.jpg}
% \videoframe{figures/video_frames/train_original_00011.jpg}
% \videoframe{figures/video_frames/train_original_00017.jpg}
% %\captionof{figure}{Example of directed camera movement }
% %\end{figure}
% {\textbf{Original Video}
% \ \\
% \ \\}
% \end{subfigure}


% \begin{subfigure}[b]{1.0\linewidth}
% \centering
% \videoframe{figures/video_frames/train_outpaint_00001.jpg}
% \videoframe{figures/video_frames/train_outpaint_00006.jpg}
% \videoframe{figures/video_frames/train_outpaint_00011.jpg}
% \videoframe{figures/video_frames/train_outpaint_00017.jpg}
% {\textbf{Outpainted Video}
% \ \\
% \ \\}
% \end{subfigure}


% \begin{subfigure}[b]{1.0\linewidth}
% \centering
% \videoframe{figures/video_frames/train_food_00001.jpg}
% \videoframe{figures/video_frames/train_food_00006.jpg}
% \videoframe{figures/video_frames/train_food_00011.jpg}
% \videoframe{figures/video_frames/train_food_00017.jpg}
% {\textbf{Stylized Video} \\
% \textbf{Prompt:} A gingerbread and candy train on a track
% \ \\}
% \end{subfigure}

% \caption{\textbf{Example of zero-shot video editing via task chaining (outpainting and stylization)} -- the original video is first outpainted and then stylized via a text prompt.\lu{We can see a clear artifact in the middle row over the train (white block). This image may raise question.}}
% \label{fig:outpaint_stylize}
% \end{figure}
